\documentclass{situationist-preprint}

\title{A Weak Duty to Obey the Law}
\msid{POL-03-a}
\msfamily{POL-03}
\msdiscipline{Political Philosophy}
\mstype{Principal manuscript}
\mskeywords{political obligation; duty to obey the law; fair play; principle of fairness; public goods; free-riding; unjust laws}

\begin{document}
\maketitle

\begin{abstract}
Is there a duty to obey the law? Consent theories find too little consent to ground one; philosophical anarchists conclude that there is no general duty, not even a prima facie one. We argue that there is a weak, defeasible duty to obey just laws, grounded in fairness. Those who receive the indispensable goods of a reasonably just legal order, such as security and the rule of law, owe it to the others whose compliance produces those goods to bear their share of the burden of compliance. The duty is weak because the wrong of a single breach is the modest wrong of free-riding, not the wrong of causing harm, so that reasons of moderate strength outweigh it. It is defeasible because it holds only while its conditions hold: the scheme must be reasonably just, its burdens fairly distributed, and the law in question not seriously unjust. For seriously unjust laws the duty lapses, since compliance with them is not a fair share of a fair burden. We answer Nozick's public address system, Simmons's objection from non-voluntary receipt, and Smith's argument that harmless breaches wrong no one.
\end{abstract}

\section{The problem}

In the \emph{Crito}, Socrates argues that by remaining in Athens he had agreed to abide by its laws, and that to escape would be to injure the city to which he owed his upbringing and his security \citep[50a--54d]{PlatoCrito}. The dialogue contains, in compressed form, three of the grounds that later theories of political obligation have developed: consent, gratitude and fairness. The question it raises is whether citizens have a moral duty to obey the law because it is the law, a duty that would hold even where the conduct the law requires is not independently required by morality.

Three answers are current. The first is that there is a general duty to obey, grounded in consent \citep[\S\S 95--122]{Locke1689} or in a natural duty to support just institutions \citep[\S 19]{Rawls1971}. The second is philosophical anarchism: there is no general duty to obey the law, not even a prima facie one \citep{Smith1973,Simmons1979,Raz1979}. The third, which we defend, is the fair-play account, on which the duty is owed by those who benefit from a cooperative scheme of rules to those whose compliance produces the benefit \citep{Hart1955,Rawls1964,Klosko1992}.

Our thesis is that the fair-play account succeeds, but only for a duty with three features that its defenders have not always acknowledged together. It is a duty to obey \emph{just} laws, not all laws. It is \emph{weak}: it can be outweighed by reasons of moderate strength. And it is \emph{defeasible}: it holds only while its conditions hold, and in particular it lapses for seriously unjust laws. These features are not concessions extracted by the anarchist's objections. They follow from the ground of the duty, and they are what allow it to survive the objections.

The question is distinct from that of legitimacy. POL-02 argues that a state is legitimate, that is, permitted to make and coercively enforce rules, if and only if it can be justified to every reasonable person subject to it. Legitimacy concerns the state's permission; obligation concerns the subject's duty.

\section{Framework}

\begin{definition}[Cooperative scheme]\label{d:scheme}
A \emph{cooperative scheme} is a practice in which a number of persons restrict their liberty according to rules, so that their general compliance produces goods available to the participants.
\end{definition}

\begin{definition}[Just law]\label{d:just}
A \emph{just law} is a rule of a legal system that is reasonably just as a whole, which is not itself seriously unjust, and whose general observance is part of a cooperative scheme from which those subject to it benefit.
\end{definition}

Definition~\ref{d:just} is the definition used in ETH-17-a, and the duty we defend is the duty that manuscript assumes. A law is \emph{seriously unjust} when it imposes grave burdens on some for reasons they could not reasonably accept, or requires those subject to it to participate in grave wrongs against others, as laws establishing slavery or persecuting a minority do. A law may be imperfectly just without being seriously unjust.

\begin{principle}[Fairness]\label{p:fair}
If a person receives goods produced by a cooperative scheme that is reasonably just and whose burdens are fairly distributed, and the goods are worth their cost to her, then she owes it to the other participants to bear her share of the burdens the scheme imposes.
\end{principle}

Principle~\ref{p:fair} is Hart's principle with the restrictions that \citet{Klosko1992} and \citet{Arneson1982} have shown it needs. The scheme must be reasonably just, since a scheme of oppression cannot generate obligations to sustain it. The burdens must be fairly distributed, since a participant who is asked to bear more than her share is not free-riding when she declines the excess. And the goods must be worth their cost to her, which rules out schemes that thrust upon people goods they do not want at a price they would not pay.

\begin{assumption}[The legal order as a scheme]\label{a:legal}
The legal order of a reasonably just state is a cooperative scheme in the sense of Definition~\ref{d:scheme}. Its goods include security of person and property, the rule of law, and the coordination of conduct on matters where coordination is needed. These goods are \emph{presumptively beneficial}: they are worth their cost to everyone, whatever her conception of the good.
\end{assumption}

\begin{assumption}[Marginality]\label{a:marginal}
The goods of a legal order depend on general, not universal, compliance. A single breach of a just law typically makes a negligible difference to the supply of those goods.
\end{assumption}

Assumption~\ref{a:legal} is Klosko's: the goods of public order are needed by everyone to pursue any plan of life. Assumption~\ref{a:marginal} is a familiar fact about public goods, on which the anarchist relies: if my breach makes no difference, what wrong does it do?

\section{The argument}

\begin{lemma}[Existence]\label{l:exist}
Those subject to a reasonably just legal order owe it to its other participants to comply with its just laws.
\end{lemma}

\begin{proof}
By Assumption~\ref{a:legal}, the legal order of a reasonably just state is a cooperative scheme whose goods are worth their cost to each person subject to it. Each such person receives those goods, since they are available to everyone within the territory and cannot be declined. By Definition~\ref{d:just}, the burdens just laws impose are part of a reasonably just scheme. By Principle~\ref{p:fair}, each person owes it to the other participants to bear her share of those burdens, and her share of the burden of a law is compliance with it.
\end{proof}

\begin{lemma}[Weakness]\label{l:weak}
The duty of Lemma~\ref{l:exist} is weak: the wrong of a single breach is the wrong of taking a share of the scheme's goods without bearing a share of its burdens, and reasons of moderate strength outweigh it.
\end{lemma}

\begin{proof}
By Assumption~\ref{a:marginal}, a single breach typically makes a negligible difference to the goods of the scheme. Its wrong is therefore not that it harms the other participants, but that it is unfair to them: the breaching person benefits from their compliance while exempting herself from the burden that produces the benefit. The weight of this wrong is the weight of the unfair advantage she takes, which is the size of her share of the burden. For most laws on most occasions, that share is small: the cost of waiting at a red light on an empty road, of paying one's share of a tax, of filing a form. A reason against the breach whose weight is the size of a small share of a burden is outweighed by any reason for the breach that is weightier than that share, which includes reasons of moderate strength such as the prevention of substantial harm to others.
\end{proof}

The weakness of the duty is fixed by its ground. A duty grounded in harm would have the weight of the harm, and since a single breach typically causes none, it would have no weight; that is Smith's argument. A duty grounded in fairness has the weight of the unfair advantage, which is real but modest.

\begin{lemma}[Defeasibility]\label{l:defeat}
The duty of Lemma~\ref{l:exist} holds only while the conditions of Principle~\ref{p:fair} hold. In particular, it lapses for seriously unjust laws.
\end{lemma}

\begin{proof}
The duty is derived from Principle~\ref{p:fair}, whose antecedent requires that the scheme be reasonably just and its burdens fairly distributed. Where these conditions fail, the principle does not apply, and no other premise of the argument yields the duty. Consider a seriously unjust law. The burden it imposes is not a fair share of a fair burden: either it falls on its victims for reasons they could not reasonably accept, in which case their compliance is not a share of a fairly distributed burden, or it requires others to participate in a grave wrong, in which case compliance is not a burden that fairness to the other participants can require, since no one can owe it to others, as a matter of fairness, to join them in wronging a third party. In neither case does the person who breaks the law take for herself a benefit others pay for while declining her share. The fairness ground is absent, and the duty lapses with respect to that law, even when the legal order as a whole remains reasonably just.
\end{proof}

\begin{theorem}\label{t:main}
There is a weak, defeasible duty to obey just laws, grounded in fairness, and it lapses for seriously unjust laws.
\end{theorem}

\begin{proof}
By Lemmas~\ref{l:exist}, \ref{l:weak} and~\ref{l:defeat}.
\end{proof}

Two features of the result deserve comment. First, the duty is owed to the other participants in the scheme, not to the state, which is the scheme's organizing agent rather than the creditor of the obligation; this is why the duty attaches to a just law whatever its content. Second, the lapse is law-specific. A legal order containing a seriously unjust law may be reasonably just as a whole: the fugitive slave laws of the antebellum United States released citizens from the duty to obey those laws, not from the duty to obey the laws of contract.

\section{Objections}

\begin{objection}[The public address system]
\citet[ch.~5]{Nozick1974} imagines neighbours who institute a public address system on which each in turn broadcasts for a day. One has enjoyed the broadcasts. Is one bound to take one's turn? It seems not, and if not, receiving benefits from a scheme does not bind.
\end{objection}

\begin{reply}
The case fails the third condition of Principle~\ref{p:fair}: the broadcasts are not worth their cost to the person, who may prefer a day of his own to a year of others' entertainment. Assumption~\ref{a:legal} concerns goods of a different kind. Security and the rule of law are worth their cost to everyone, since without them no plan of life can be pursued with confidence. Nozick's case shows that Hart's unrestricted principle is too strong, and the restricted principle concedes this; it does not show that the restricted principle fails for presumptive public goods \citep{Klosko1992}.
\end{reply}

\begin{objection}[Non-voluntary receipt]
\citet{Simmons1979} argues that the fairness principle binds only those who have \emph{accepted} the benefits of a scheme, in the sense of taking them willingly and knowingly, and not those who merely receive them. Most citizens merely receive the goods of the legal order, which they cannot avoid. So the principle cannot ground a general duty.
\end{objection}

\begin{reply}
Acceptance matters where benefits are optional, because there the free choice to take a benefit is what makes it unfair to decline the cost. It does not matter where the benefits are indispensable and worth their cost to everyone, because there the unfairness of declining the cost does not depend on a choice to take the benefit. A person who would accept security if offered it on fair terms, and who receives it without being able to decline it, takes the same advantage over those who pay for it by refusing to pay as one who chose it. Requiring acceptance in such cases would give a premium to those who cannot avoid a good over those who choose it, which is the opposite of what fairness requires \citep{Arneson1982}.
\end{reply}

\begin{objection}[Harmless breaches]
\citet{Smith1973} argues that there is no prima facie obligation to obey the law, since a prima facie obligation must have some weight in every case it covers, and many breaches, such as running a red light on an empty road at night, have no bad consequences and wrong no one.
\end{objection}

\begin{reply}
Smith's argument assumes that a breach wrongs someone only if it harms someone. The fairness ground denies this. The driver who runs the red light on the empty road takes a liberty that others forgo, and the system of traffic lights works because they forgo it; he benefits from their compliance, since it makes the roads safe for him on other nights, while declining to comply. The wrong is small, which is why the duty is weak, and it is a wrong, which is why the duty exists. Smith is right that the duty has little weight in such cases and that a modest reason, such as getting an injured person to hospital, outweighs it. He is wrong that it has none.
\end{reply}

\section{Scope and consequences}

The theorem does not establish that the duty to obey is the most important reason to comply with the law; most laws prohibit conduct that is independently wrong. It does not establish a duty to obey the laws of a seriously unjust regime, where the conditions of Principle~\ref{p:fair} fail. And since the duty is weak, it does not establish that a person who breaks a just law acts wrongly all things considered.

The last point is developed in ETH-17-a, which takes the present result as its premise: a breach of a just law to prevent harm is permissible when the harm is grave, the breach minimal and public, and the legal consequences accepted, since gravity and minimality outweigh a weak duty and publicity and acceptance keep the breach from being free-riding. ETH-17-a also relies on Lemma~\ref{l:defeat}: breaches of seriously unjust laws are not governed by its conditions, since the duty does not extend to such laws. ETH-26-a argues that punishment communicates a censure the offender has reason to accept, and the fairness duty is one such reason. POL-23 argues that revolution is justified under conditions of extreme injustice, in which the scheme is no longer reasonably just and the duty lapses for the legal order as a whole.

POL-11 argues that the social contract exists as an implicit framework of mutual expectations. That framework is the cooperative scheme of Definition~\ref{d:scheme} viewed as a social fact, and the present result supplies part of its normative force where its terms are reasonably just.

\section{Conclusion}

There is a duty to obey the law, but it is a duty to obey just laws, it is weak, and it is defeasible. Each feature follows from its ground in fairness. Those who receive the indispensable goods of a reasonably just legal order owe it to those whose compliance produces them to bear their share; the wrong of declining is the modest wrong of free-riding, not the wrong of harm; and where a law is seriously unjust, compliance is no fair share of a fair burden, and the duty lapses. The anarchist is right that there is no strong, general duty to obey whatever the state commands. The fair-play theorist is right that there is a duty all the same.

\begin{thebibliography}{99}\small
\bibitem[Arneson(1982)]{Arneson1982} Arneson, R.~J. (1982). The principle of fairness and free-rider problems. \emph{Ethics}, 92(4), 616--633.
\bibitem[Hart(1955)]{Hart1955} Hart, H.~L.~A. (1955). Are there any natural rights? \emph{The Philosophical Review}, 64(2), 175--191.
\bibitem[Klosko(1992)]{Klosko1992} Klosko, G. (1992). \emph{The Principle of Fairness and Political Obligation}. Rowman \& Littlefield.
\bibitem[Locke(1689)]{Locke1689} Locke, J. (1689). \emph{Two Treatises of Government}. Cited by section of the Second Treatise.
\bibitem[Nozick(1974)]{Nozick1974} Nozick, R. (1974). \emph{Anarchy, State, and Utopia}. Basic Books.
\bibitem[Plato(Crito)]{PlatoCrito} Plato. \emph{Crito}. Cited by Stephanus pages.
\bibitem[Rawls(1964)]{Rawls1964} Rawls, J. (1964). Legal obligation and the duty of fair play. In S.~Hook (ed.), \emph{Law and Philosophy}. New York University Press.
\bibitem[Rawls(1971)]{Rawls1971} Rawls, J. (1971). \emph{A Theory of Justice}. Harvard University Press.
\bibitem[Raz(1979)]{Raz1979} Raz, J. (1979). \emph{The Authority of Law: Essays on Law and Morality}. Clarendon Press.
\bibitem[Simmons(1979)]{Simmons1979} Simmons, A.~J. (1979). \emph{Moral Principles and Political Obligations}. Princeton University Press.
\bibitem[Smith(1973)]{Smith1973} Smith, M.~B.~E. (1973). Is there a prima facie obligation to obey the law? \emph{The Yale Law Journal}, 82(5), 950--976.
\end{thebibliography}

\end{document}
